%-------------------------
% Resume in Latex
% Author : Harshibar
% Based off of: https://github.com/jakeryang/resume
% License : MIT
%------------------------

\documentclass[letterpaper,10pt]{article}

\usepackage{latexsym}
\usepackage[empty]{fullpage}
\usepackage{titlesec}
\usepackage{marvosym}
\usepackage[usenames,dvipsnames]{color}
\usepackage{verbatim}
\usepackage{enumitem}
\usepackage[hidelinks]{hyperref}
\usepackage{fancyhdr}
\usepackage[T2A]{fontenc}
\usepackage[utf8]{inputenc}
\usepackage[english,russian]{babel}
\usepackage{tabularx}
% only for pdflatex
% \input{glyphtounicode}

% fontawesome
\usepackage{fontawesome5}

% fixed width
\usepackage[scale=0.90,lf]{FiraMono}

% light-grey
\definecolor{light-grey}{gray}{0.83}
\definecolor{dark-grey}{gray}{0.3}
\definecolor{text-grey}{gray}{.08}

\DeclareRobustCommand{\ebseries}{\fontseries{eb}\selectfont}
\DeclareTextFontCommand{\texteb}{\ebseries}

% custom underilne
\usepackage{contour}
\usepackage[normalem]{ulem}
\renewcommand{\ULdepth}{1.8pt}
\contourlength{0.8pt}
\newcommand{\myuline}[1]{%
  \uline{\phantom{#1}}%
  \llap{\contour{white}{#1}}%
}
\newcommand{\latinmono}[1]{{\fontencoding{T1}\ttfamily #1}}


% custom font: helvetica-style with Cyrillic support
\usepackage[default,scale=0.92]{opensans}


\pagestyle{fancy}
\fancyhf{} % clear all header and footer fields
\fancyfoot{}
\setlength{\footskip}{4.1pt}
\renewcommand{\headrulewidth}{0pt}
\renewcommand{\footrulewidth}{0pt}

% Adjust margins
\addtolength{\oddsidemargin}{-0.5in}
\addtolength{\evensidemargin}{0in}
\addtolength{\textwidth}{1in}
\addtolength{\topmargin}{-.5in}
\addtolength{\textheight}{1.0in}

\urlstyle{same}

\raggedbottom
\raggedright
\setlength{\tabcolsep}{0in}

% Sections formatting - serif
% \titleformat{\section}{
%   \vspace{2pt} \scshape \raggedright\large % header section
% }{}{0em}{}[\color{black} \titlerule \vspace{-5pt}]

% TODO EBSERIES
% sans serif sections
\titleformat {\section}{
    \bfseries \vspace{0pt} \raggedright \large % header section
}{}{0em}{}[\color{light-grey} {\titlerule[2pt]} \vspace{-6pt}]

% only for pdflatex
% Ensure that generate pdf is machine readable/ATS parsable
% \pdfgentounicode=1

%-------------------------
% Custom commands
\newcommand{\resumeItem}[1]{
  \item\small{
    {#1 \vspace{-2pt}}
  }
}

\newcommand{\resumeSubheading}[4]{
  \vspace{-1pt}\item
    \begin{tabular*}{\textwidth}[t]{l@{\extracolsep{\fill}}r}
      \textbf{#1} & {\color{dark-grey}\small #2}\vspace{1pt}\\ % top row of resume entry
      \textit{#3} & {\color{dark-grey} \small #4}\\ % second row of resume entry
    \end{tabular*}\vspace{-4pt}
}

\newcommand{\resumeSubSubheading}[2]{
    \item
    \begin{tabular*}{\textwidth}{l@{\extracolsep{\fill}}r}
      \textit{\small#1} & \textit{\small #2} \\
    \end{tabular*}\vspace{-7pt}
}

\newcommand{\resumeProjectHeading}[2]{
    \item
    \begin{tabular*}{\textwidth}{l@{\extracolsep{\fill}}r}
      #1 & {\color{dark-grey}\small #2} \\
    \end{tabular*}\vspace{-4pt}
}

\newcommand{\resumeSubItem}[1]{\resumeItem{#1}\vspace{-4pt}}

\renewcommand\labelitemii{$\vcenter{\hbox{\tiny$\bullet$}}$}

% CHANGED default leftmargin  0.15 in
\newcommand{\resumeSubHeadingListStart}{\begin{itemize}[leftmargin=0in, label={}, itemsep=5pt, topsep=0pt]}
\newcommand{\resumeSubHeadingListEnd}{\end{itemize}}
\newcommand{\resumeItemListStart}{\begin{itemize}[leftmargin=0.18in, itemsep=0pt, topsep=1pt, parsep=0pt, partopsep=0pt]}
\newcommand{\resumeItemListEnd}{\end{itemize}\vspace{-2pt}}

\color{text-grey}

%-------------------------------------------
%%%%%%  RESUME STARTS HERE  %%%%%%%%%%%%%%%%%%%%%%%%%%%%


\begin{document}

%----------HEADING----------
\begin{center}
    \textbf{\Large Очкин Никита Валерьевич} \\ \vspace{5pt}
    \small \faPhone* \latinmono{+7 977 194-80-50} \hspace{1pt} $|$
    \hspace{1pt} \faEnvelope \hspace{2pt} \latinmono{nkt.ochkin@gmail.com} \hspace{1pt} $|$
    \hspace{1pt} \faTelegram \hspace{2pt} 
    \href{https://t.me/namenick9}{\myuline{\latinmono{@namenick9}}} \hspace{1pt} $|$
    \hspace{1pt} \faGithub \hspace{2pt}
    \href{https://github.com/namenick91}{\myuline{\latinmono{github.com/namenick91}}} \hspace{1pt} $|$
    \hspace{1pt} \faMapMarker* \hspace{2pt} Москва
    \\ \vspace{-3pt}
\end{center}


%-----------EDUCATION-----------
\section{ОБРАЗОВАНИЕ}
  \resumeSubHeadingListStart
    \resumeSubheading
      {МГТУ им. Н.Э. Баумана, Факультет <<Фундаментальные науки>>}{2022 -- 2026}
      {Бакалавриат, 02.03.01 <<Математика и компьютерные науки>>}{}
      	\resumeItemListStart
    	\resumeItem{\textbf{Профиль:} суперкомпьютерное моделирование и искусственный интеллект в инженерных задачах.}
      \resumeItem{\textbf{ВКР:} <<Разработка алгоритма вопросно-ответной системы по курсу высшей математики на базе большой языковой модели>>.}
      \resumeItemListEnd
  \resumeSubHeadingListEnd

\vspace{-7pt}

%-----------EXPERIENCE-----------
\section{ОПЫТ РАБОТЫ}
  \resumeSubHeadingListStart

    \resumeSubheading
      {Институт системного программирования им. В.\,П.~Иванникова РАН}{янв. 2026 -- н. в.}
      {Лаборант, 16-й отдел информационных систем}{}
      \resumeItemListStart
        \resumeItem{Провёл сравнительное исследование VLM OCR-систем на различных типах документов (6 систем $\times$ 5 сценариев): инференс на GPU A100, метрики для каждого сценария (CER/WER, PER, CDM).}
        \resumeItem{Разработал FastAPI-сервис для парсинга документов на базе PaddleOCR-VL-1.5: vLLM-инференс, Docker-развёртывание; результат OCR преобразуется в структурированный граф документа. Micro-CER сервиса --- 3,7 \%.}
        \resumeItem{Разработал бенчмарк для выбора VLM, извлекающей табличные данные из разнотипных графиков (chart-to-table), с целью дальнейшего внедрения в сервис: 11 моделей, 5 датасетов, 4 метрики (8 отчётных конфигураций) с bootstrap-CI.}
        \resumeItem{Рецензировал научные статьи по анализу документов методами DL.}
      \resumeItemListEnd

    \resumeSubheading
      {МГТУ им. Н.Э. Баумана}{окт. 2024 -- н. в.}
      {Техник, кафедра ФН-11 <<Вычислительная математика и математическая физика>>}{}
      \resumeItemListStart
        \resumeItem{Разработал вопросно-ответную RAG-систему для образовательной платформы NOMOTEX: ETL разнородных источников, гибридный поиск (BGE-M3) и хранение (Qdrant + PostgreSQL), диалоговая память, ссылки на источники в ответах модели.}
        \resumeItem{Оценил качество (RAGAS) и выбрал конфигурацию системы. Сейчас --- внедряю на платформу.}
      \resumeItemListEnd

  \resumeSubHeadingListEnd

\vspace{-7pt}

%-----------RESEARCH-----------
\section{НАУЧНАЯ РАБОТА}
    \resumeSubHeadingListStart
      \resumeProjectHeading
          {\textbf{Convformer}: прогнозирование временных рядов}{2025 -- 2026}
          \resumeItemListStart
          \resumeItem{Расширил эмбеддинг Informer двухслойным свёрточным блоком, добавил автоматическую декомпозицию ряда из Autoformer и заменил ProbSparse на механизм внимания FAVOR+ (Performer), работающий за линейное время.}
          \resumeItem{На длинных горизонтах получил меньший MSE на стандартных бенчмарках (ETT, ECL и др.).}
            \resumeItem{Препринт статьи готов, идёт работа над публикацией в журнале уровня Q2--Q3; код и препринт: \href{https://github.com/namenick91/Convformer}{\myuline{GitHub}}.}
          \resumeItemListEnd

    \resumeSubHeadingListEnd

\vspace{-7pt}

%-----------HACKATHONS-----------
\section{ХАКАТОНЫ}
    \resumeSubHeadingListStart
      \resumeProjectHeading
          {\textbf{МТС True Tech Hack 2026}: GPTHub}{2026}
          \resumeItemListStart
            \resumeItem{Возглавлял команду из 5 человек.}
            \resumeItem{Собрали платформу на базе OpenWebUI: маршрутизация запросов, RAG, долгосрочная память на PostgreSQL + pgvector и Deep Research через gpt-researcher.}
            \resumeItem{Отвечал за архитектуру RAG, подбор моделей и пресетов, двухуровневую память, интеграцию Deep Research и координацию команды.}
          \resumeItemListEnd
    \resumeSubHeadingListEnd

\vspace{-7pt}

%-----------PROJECTS-----------
\section{PET-ПРОЕКТЫ}
    \resumeSubHeadingListStart
      \resumeProjectHeading
          {\textbf{Customer Churn} (CatBoost)}{\href{https://github.com/namenick91/ML-Workshelf\#customer-churn-prediction}{\myuline{GitHub}}}
          \resumeItemListStart
            \resumeItem{Построил модель оттока клиентов: очистка и типизация данных, генерация признаков, k-fold и подбор гиперпараметров в Optuna.}
          \resumeItemListEnd

      \resumeProjectHeading
          {\textbf{Multimodal RAG} для поиска фильмов}{\href{https://github.com/namenick91/rag-multimodal-movies}{\myuline{GitHub}}}
          \resumeItemListStart
            \resumeItem{Собрал сервис: FastAPI + Qdrant + MinIO + Streamlit; использовал E5/CLIP-эмбеддинги и локальную модель Qwen для суммаризации.}
          \resumeItemListEnd

      \resumeProjectHeading
          {\textbf{Simpsons Character Classification} (EfficientNetV2-S)}{\href{https://github.com/namenick91/ML-Workshelf\#simpsons-character-classification}{\myuline{GitHub}}}
          \resumeItemListStart
            \resumeItem{Обучил многоклассовый классификатор изображений: двухфазный fine-tuning, RandAugment, CutMix, MixUp, AdamW + Cosine и early stopping.}
          \resumeItemListEnd

      \resumeProjectHeading
          {\textbf{Semantic Segmentation} (U-Net / SegNet)}{\href{https://github.com/namenick91/ML-Workshelf\#image-segmentation-with-cnns}{\myuline{GitHub}}}
          \resumeItemListStart
            \resumeItem{Сравнил модели для бинарной сегментации медицинских изображений; обучал с функциями потерь BCE/Dice/Focal, качество оценивал по IoU, применял early stopping по валидационной метрике.}
          \resumeItemListEnd
    \resumeSubHeadingListEnd

\vspace{-7pt}

%-----------PROGRAMMING SKILLS-----------
\section{НАВЫКИ}
 \begin{itemize}[leftmargin=0in, label={}]
    \small{\item{
     \textbf{Инструменты:}{ Python, SQL, NumPy, SciPy, SymPy, pandas, FastAPI, Streamlit, \LaTeX}\vspace{2pt} \\
     \textbf{ML / DL:}{ PyTorch, scikit-learn, Hugging Face Transformers, RAG; оценка и бенчмаркинг LLM/VLM, прогнозирование временных рядов}\vspace{2pt} \\
     \textbf{Инфраструктура:}{ Git, Linux, bash, Docker, vLLM, pytest, PostgreSQL + pgvector, Qdrant, MinIO}\vspace{2pt} \\
     \textbf{Языки:}{ русский --- родной; английский --- B2--C1}
    }}
 \end{itemize}


%-------------------------------------------
\end{document}
